%% Chapter 7 -- Backend Architecture (RPi Bridge)
\chapter{Backend Architecture}
\label{ch:backend}

\section{Raspberry Pi Role}

The Raspberry Pi is the communication hub between the local MQTT network and Firebase cloud. It runs two services:
\begin{itemize}[leftmargin=1.5em]
    \item \textbf{Mosquitto}: MQTT message broker; all ESP32 nodes connect to it
    \item \textbf{bridge.py}: Python service that subscribes to MQTT topics and synchronises data with Firebase
\end{itemize}

The RPi is connected via Ethernet to the lab router and must have a static IP address (192.168.0.10 by default).

\section{Mosquitto Configuration}

Mosquitto is configured in \texttt{/etc/mosquitto/conf.d/mms.conf}:

\begin{lstlisting}[caption={Mosquitto configuration file}]
listener 1883 0.0.0.0
allow_anonymous true          # change to false after adding auth
persistence true
persistence_location /var/lib/mosquitto/
log_dest syslog
log_type all
log_timestamp true
max_connections 50
max_queued_messages 1000
\end{lstlisting}

Key settings:
\begin{itemize}[leftmargin=1.5em]
    \item \textbf{Persistence}: Retained messages (config push, revocation list, node status) survive broker restarts. Nodes reconnecting after a broker restart immediately receive their current config and revocation list.
    \item \textbf{max\_connections}: 50 is sufficient for 15 nodes + 1 bridge + 5 admin tools.
    \item \textbf{allow\_anonymous}: Set to \texttt{true} during initial deployment. Should be changed to \texttt{false} with password authentication after all nodes are confirmed working. See Section~\ref{sec:mqtt_auth}.
\end{itemize}

\section{bridge.py Architecture}

\texttt{bridge.py} is a single-process Python script that runs the following listeners concurrently using Paho MQTT's event loop and Firebase SDK callbacks:

\begin{enumerate}[leftmargin=1.5em]
    \item \textbf{MQTT subscriber: \texttt{mms/nodes/+/event}} $\rightarrow$ write session to Firestore
    \item \textbf{MQTT subscriber: \texttt{mms/nodes/+/status}} $\rightarrow$ update RTDB node state
    \item \textbf{RTDB listener: \texttt{/mms/commands/\{id\}}} $\rightarrow$ forward commands to MQTT
    \item \textbf{RTDB listener: \texttt{/mms/ota}} $\rightarrow$ broadcast OTA command to all nodes
    \item \textbf{Firestore listener: \texttt{/machines/*}} $\rightarrow$ push config to nodes via retained MQTT
    \item \textbf{Firestore listener: \texttt{/revoked}} $\rightarrow$ push revocation list to \texttt{mms/all/revoked}
    \item \textbf{Periodic sweep (every 10 min)} $\rightarrow$ clean up stale unacknowledged RTDB commands
\end{enumerate}

\section{Session Event Processing}

When bridge.py receives an event on \texttt{mms/nodes/\{id\}/event}:

\begin{lstlisting}[language=Python,caption={Session event handler in bridge.py (simplified)}]
def on_event(client, userdata, message):
    event = json.loads(message.payload)
    machine_id = extract_machine_id(message.topic)

    # Deduplication: check for existing session within +/-5s
    start_ts = event["t"] - event["d"]
    existing = db.collection("sessions") \
        .where("machine_id", "==", str(machine_id)) \
        .where("start_time", ">=", datetime.fromtimestamp(start_ts - 5)) \
        .where("start_time", "<=", datetime.fromtimestamp(start_ts + 5)) \
        .limit(1).get()

    if existing:
        logger.info(f"Duplicate session detected -- skipping")
        return

    # Get user info from Firestore
    user_doc = db.collection("users").document(event["r"]).get()
    user_name = user_doc.get("name") if user_doc.exists else "Unknown"

    # Write session record
    session = {
        "machine_id":   str(machine_id),
        "roll_number":  event["r"],
        "user_name":    user_name,
        "start_time":   datetime.fromtimestamp(event["t"] - event["d"]),
        "end_time":     datetime.fromtimestamp(event["t"]),
        "duration_s":   event["d"],
        "end_reason":   event["er"],
        "status":       "completed",
    }
    db.collection("sessions").add(session)

    # Update user last_seen
    db.collection("users").document(event["r"]).update(
        {"last_seen": firestore.SERVER_TIMESTAMP}
    )
\end{lstlisting}

\section{Command Forwarding}

When bridge.py's RTDB listener fires for a command:

\begin{lstlisting}[language=Python,caption={Command forward handler in bridge.py (simplified)}]
def on_commands_change(event):
    machine_id = extract_id_from_path(event.path)
    data = event.data

    if data is None or "command" not in data:
        return

    command = data["command"]
    ts = data.get("timestamp", 0)

    # Stale command check (5-minute TTL)
    if time.time() - ts > 300:
        logger.info(f"Stale command for machine {machine_id} -- discarding")
        rtdb.reference(f"mms/commands/{machine_id}").update({
            "acknowledged": True,
            "discarded": True,
            "reason": "stale_command",
        })
        return

    # Forward to node via MQTT
    payload = json.dumps({"cmd": command, "issued_by": data.get("issued_by"), "ts": ts})
    mqtt_client.publish(f"mms/nodes/{machine_id}/command", payload, qos=1)
\end{lstlisting}

\section{Revocation List Management}

The revocation list is maintained in Firestore \texttt{/revoked} collection. When any document is added or removed:

\begin{lstlisting}[language=Python,caption={Revocation push in bridge.py (simplified)}]
def on_revoked_change(col_snapshot, changes, read_time):
    # Rebuild full list from snapshot
    uids = [doc.id for doc in col_snapshot]

    payload = json.dumps({
        "uids": uids,
        "updated_at": int(time.time())
    })

    # Publish retained -- nodes receive it on connect even if offline during update
    mqtt_client.publish("mms/all/revoked", payload, qos=1, retain=True)
    logger.info(f"Revocation list pushed: {len(uids)} UID(s)")
\end{lstlisting}

\section{Systemd Service}

bridge.py runs as a systemd service with automatic restart:

\begin{lstlisting}[caption={/etc/systemd/system/mms-bridge.service}]
[Unit]
Description=MMS Firebase Bridge
After=network-online.target mosquitto.service
Wants=network-online.target

[Service]
Type=notify
User=pi
WorkingDirectory=/home/pi/mms
EnvironmentFile=/home/pi/mms/.env
ExecStart=/home/pi/mms/venv/bin/python3 /home/pi/mms/bridge.py
Restart=always
RestartSec=5
StartLimitIntervalSec=60
StartLimitBurst=5
WatchdogSec=60
NotifyAccess=main

[Install]
WantedBy=multi-user.target
\end{lstlisting}

Key settings:
\begin{itemize}[leftmargin=1.5em]
    \item \texttt{Restart=always}: restart on any exit (crash, OOM, signal)
    \item \texttt{StartLimitBurst=5}: stop restarting after 5 failures in 60s (prevents runaway restart loop)
    \item \texttt{WatchdogSec=60}: systemd kills and restarts the bridge if it stops sending watchdog pings within 60s
    \item \texttt{EnvironmentFile}: loads \texttt{.env} before starting (MQTT credentials, Firebase URL)
\end{itemize}

\section{Health Check Endpoint}

bridge.py serves a minimal HTTP health check on port 8080:

\begin{lstlisting}[style=bashstyle]
curl http://192.168.0.10:8080/health
# Returns:
{"mqtt_connected": true, "fb_connected": true, "uptime_s": 3600}
\end{lstlisting}

This endpoint can be polled by a monitoring system or checked manually to confirm the bridge is alive.

\section{MQTT Authentication}
\label{sec:mqtt_auth}

In the initial deployment, \texttt{allow\_anonymous true} is used to simplify setup. Once all nodes are confirmed working, authentication should be enabled:

\begin{lstlisting}[style=bashstyle]
# Create password file
sudo mosquitto_passwd -c /etc/mosquitto/passwd bridge
sudo mosquitto_passwd /etc/mosquitto/passwd node_1
# ... repeat for each node (node_1 through node_5)

# Update mms.conf:
# Change: allow_anonymous true
# To:     allow_anonymous false
#         password_file /etc/mosquitto/passwd

sudo systemctl restart mosquitto
\end{lstlisting}

Update each node's \texttt{secrets.h} with \texttt{MQTT\_USER} and \texttt{MQTT\_PASSWORD} and re-flash.

For per-node topic ACLs (highly recommended before production), add an ACL file:

\begin{lstlisting}[caption={/etc/mosquitto/acl (per-node topic restrictions)}]
# Bridge: full access
user bridge
topic readwrite mms/#

# Per-node: scoped to own topics + global revocation
user node_1
topic write mms/nodes/1/status
topic write mms/nodes/1/event
topic write mms/nodes/1/heartbeat
topic read  mms/nodes/1/command
topic read  mms/nodes/1/config
topic read  mms/all/revoked
\end{lstlisting}

Without ACLs, a compromised node could issue commands to other nodes or trigger OTA for the entire lab.
